\documentclass[11pt,oneside,openany]{book}

% 東吳大學法研所「憲法專題研究（一）」第 11 週（2026-11-23）預習問答集：
% Michelman, Brennan and Democracy, 86 Calif. L. Rev. 399 (1998)；Post, Democracy, Popular Sovereignty, and Judicial Review, 86 Calif. L. Rev. 429 (1998)。
% 版式照 ../W10_預習/W10_Law_預習.tex。編譯：./build.sh。

\def\TransLaTeXDir{../../../../../1142/Translation/LaTeX}
\input{\TransLaTeXDir/course-notes.sty}
\addbibresource{refs.bib}
\CourseNotesMeta{民主理論與違憲審查：第 11 週預習}{王逸帆}
\CourseNotesChinese[部]
\CourseNotesLawIndexes[noautomatic]

\providecommand{\tightlist}{\setlength{\itemsep}{0pt}\setlength{\parskip}{0pt}}
\makeatletter\def\LTcaptype{table}\makeatother

\newcommand{\pg}[1]{（#1）}
\newcommand{\en}[1]{\textit{#1}}
\newif\ifshowindex \showindextrue
\newcommand{\ix}[2]{\ifshowindex\index[#1]{#2}\fi}
\setlength{\emergencystretch}{2em}
\newcolumntype{L}[1]{>{\raggedright\arraybackslash}p{#1}}

\newcommand{\documentversion}{v0.1（2026-11-23）}

\begin{document}

\frontmatter
\coursetitle{民主理論與違憲審查}{第 11 週預習：Michelman 與 Post}%
  {東吳大學法律學系研究所　憲法專題研究（一）\\115 學年度第 1 學期　第 11 週（2026 年 11 月 23 日）}%
  {\documentversion}
\tableofcontents

\chapter*{凡例}
\addcontentsline{toc}{chapter}{凡例}
本篇為課前預習，兩篇頁碼依刊本（86 Calif. L. Rev. 399--427、429--443）。Michelman 一文所用 PDF 為 Harvard DASH 本，PDF 頁減 396 為刊本頁；Post 一文所用 PDF 為 JSTOR 本，PDF 頁減 427 為刊本頁。引文照錄英文原句，中譯附於其後的方括號內。
% 頁碼換算以 Post 引 Michelman 之頁（419、420、424、425--426、427）回核 DASH 本，皆合。Post 註 6 引 Michelman「at 430」，所引定義實在 400 頁。


\mainmatter
% 各部待大綱定案後加入
\nocite{michelman1998,post1998}
\backmatter
\printbibliography[heading=bibintoc]
\printindex[terms]
\printindex[statutes]
\printindex[cases]
\printindex[persons]

\end{document}
